\section{Lösen von Ungleichungen}
\label{sec:fundamentals_solving_inequalities}

Ungleichungen werden grundsätzlich analog zu Gleichungen umgeformt: Auf beiden
Seiten wird dieselbe zulässige Rechenoperation durchgeführt. Eine Besonderheit
tritt bei der Multiplikation oder Division mit einer negativen Zahl auf.

\begin{DefinitionBox}{Ungleichung}{fundamentals_inequality}
  Eine \emph{Ungleichung} \EnglishTerm{inequality} vergleicht zwei Terme mit
  einem der Vergleichszeichen
  \[
    <,\qquad \leq,\qquad >,\qquad \geq.
  \]
\end{DefinitionBox}

Beispielsweise ist $\begin{array}{c} 5x+2=3 \\ [0.3em] 5x+2>3 \end{array}$ eine
$\begin{array}{l} \text{Gleichung.} \\[0.3em] \text{Ungleichung.} \end{array}$ \\

Ungleichungen lassen sich ähnlich wie Gleichungen durch Äquivalenzumformungen lösen.
\begin{ExampleBox}{Lineare Ungleichung}{fundamentals_linear_inequality_example}
  Gesucht sind alle reellen Lösungen $x\in\R$ der Ungleichung $5x+2>3$.
  \begin{ClozeArea}[25]{7}
    \[
      \begin{array}{r@{\quad}l@{\qquad}l}
        & 5x+2>3 & \TransformOp{-2} \\[0.7em]
        \Longleftrightarrow & 5x>1 & \TransformOp{:5} \\[0.7em]
        \Longleftrightarrow & x>\dfrac15. &
      \end{array}
    \]
  \end{ClozeArea}
  Damit ist die Lösungsmenge
  \[
    \mathcal{L}=\Cloze{7cm}{\left(\frac15,\infty\right)}.
  \]
\end{ExampleBox}

Ein wesentlicher Unterschied zu Gleichungen besteht darin, dass sich das
Vergleichszeichen umkehren kann.
\begin{KeyPointBox}{Umkehrung des Vergleichszeichens}{fundamentals_reversing_inequality_sign}
  Wird eine Ungleichung auf beiden Seiten mit einer negativen Zahl multipliziert (oder der Kehrwert gebildet),
  kehrt sich das Vergleichszeichen um.
  \begin{align*}
      a<b \quad &\Longleftrightarrow \quad -a>-b, \\
      a\leq b \quad &\Longleftrightarrow \quad -a\geq-b.
  \end{align*}
\end{KeyPointBox}


\begin{ExampleBox}{Umkehrung des Vergleichszeichens in Ungleichungen}{fundamentals_reversing_inequality_sign_examples}
    \begin{enumerate}[label=(\alph*)]
      \item Vergleichszeichen bei Multiplikation mit einer negativen Zahl
        \[
          \displaystyle 4<5 \quad\TransformOp{\cdot(-1)}
          \qquad \Longleftrightarrow \qquad
          -4\Cloze{0.6cm}{>}-5.
        \]
      \item Vergleichszeichen bei Kehrwertbildung
        \[
          3\geq2 \quad\TransformOp{\frac{1}{\left(\ldots\right)}}
          \qquad \Longleftrightarrow \qquad
          \frac{1}{3}\Cloze{0.6cm}{\leq}\frac{1}{2}
        \]

      \item Lösen der Ungleichung $-2x+3<7$.
        \IfClozeTF{%
          \begin{ClozeArea}[31]{8}
          \end{ClozeArea}
        }{%
          \[
            \begin{array}{r@{\quad}l@{\qquad}l}
              & -2x+3<7 & \TransformOp{-3} \\[0.7em]
              \Longleftrightarrow & -2x<4
                & \TransformOp{:(-2)\quad\text{(Vergleichszeichen wird umgekehrt)}} \\[0.7em]
              \Longleftrightarrow & x>-2. &
            \end{array}
          \]
        }
        Somit ist die Lösungsmenge
        $\mathcal{L}=\Cloze{2.5cm}{(-2,\infty)}$.
    \end{enumerate}
\end{ExampleBox}
